\clearpage
\SourceLeaf{66}
\pdfbookmark[1]{Disquisitio de elementis ellipticis Palladis. Review of 1810}{anzeige-palladis1810}
\begin{center}\rule{.75\textwidth}{.3pt}\\
Göttingische gelehrte Anzeigen. 13 December 1810.\\
\rule{.75\textwidth}{.3pt}\end{center}
On 25 November Professor \textsc{Gauss} presented a paper to the Royal Society of Sciences:
\begin{center}
\textit{Disquisitio de elementis ellipticis Palladis ex oppositionibus annorum}\\
1803, 1804, 1805, 1807, 1808, 1809.
\end{center}
For three years the author had made no new calculations of the orbit of Pallas. Because the planet was so faint, the observations of 1808 had been very scanty and defective, and some of them only became known later. It therefore seemed hardly worth improving the elements on their basis at the time. Only after Professor \textsc{Gauss} had received the 1809 observations of Pallas made by \textsc{Bouvard} at the Imperial Observatory in Paris did he calculate both the opposition of 1809 and that of 1808 (see this journal, 24 February 1810). Further analysis of all six oppositions observed so far showed that an elliptic orbit no longer suffices to represent them all accurately. This is a consequence of the large perturbations that this planet experiences from the others, particularly from Jupiter. Professor \textsc{Gauss} demonstrated this in two ways. First he calculated three sets of elliptic elements, each from four oppositions, the first from those of 1803, 1804, 1805 and 1807,
\SourceLeaf{67}%
the second from those of 1804, 1805, 1807 and 1808, and the third from the oppositions of 1805, 1807, 1808 and 1809. Only small differences should have appeared among these sets if the motion were purely elliptic. We set the three sets side by side here:
\par\nobreak
\begingroup
\newcommand{\GaussVIAngle}[6]{%
  \makebox[17pt][r]{$#1$}\makebox[4pt][l]{${}#2$}\hspace{2pt}%
  \makebox[12pt][r]{$#3$}\makebox[4pt][l]{${}#4$}\hspace{2pt}%
  \makebox[12pt][r]{$#5$}\makebox[18pt][l]{${}#6$}%
}
{\centering\small
\setlength{\tabcolsep}{8pt}
\begin{tabular}{c|l|l|l}
& \multicolumn{1}{c|}{I.} & \multicolumn{1}{c|}{II.} & \multicolumn{1}{c}{III.}\\\hline
$A$ & \GaussVIAngle{221}{^0}{39}{'}{30}{''4} & \GaussVIAngle{221}{^0}{34}{'}{56}{''7} & \GaussVIAngle{221}{^0}{23}{'}{24}{''6}\\
$B$ & \multicolumn{1}{r|}{$770.2143$} & \multicolumn{1}{r|}{$770.4467$} & \multicolumn{1}{r}{$770.9265$}\\
$C$ & \GaussVIAngle{121}{}{3}{}{11}{.4} & \GaussVIAngle{121}{}{5}{}{22}{.1} & \GaussVIAngle{120}{}{58}{}{4}{.8}\\
$D$ & \GaussVIAngle{172}{}{28}{}{56}{.9} & \GaussVIAngle{172}{}{28}{}{46}{.8} & \GaussVIAngle{172}{}{27}{}{52}{.4}\\
$E$ & \GaussVIAngle{34}{}{37}{}{41}{.0} & \GaussVIAngle{34}{}{37}{}{31}{.5} & \GaussVIAngle{34}{}{36}{}{49}{.4}\\
$F$ & $0.2450198$ & $0.2447624$ & $0.2446335$\\
$G$ & $0.4423149$ & $0.4422276$ & $0.4420473$
\end{tabular}\par}
\begin{samepage}\small
where the symbols stand for the following quantities
\par\noindent\hspace*{.06\linewidth}%
\begin{tabular}{@{}l@{\enspace}p{.74\linewidth}@{}}
  $A$ & epoch of the mean longitude in 1803 for the Göttingen meridian\\
  $B$ & mean daily tropical motion\\
  $C$ & longitude of perihelion in 1803\\
  $D$ & longitude of the ascending node in 1803\\
  $E$ & orbital inclination\\
  $F$ & eccentricity\\
  $G$ & logarithm of the semi-major axis\\
\end{tabular}\par\smallskip
\end{samepage}
It will be necessary to take into account the perturbations of Pallas by the other planets if one seeks lasting agreement between appearances and calculation. For certain purposes, however, one will still prefer a purely elliptic orbit if it avoids the intolerably lengthy computations that the great number of perturbation equations would otherwise require. In particular, when the aim is merely to find the planet again, it is certainly more convenient to base purely elliptic elements on the four immediately preceding oppositions and to construct the ephemeris from them. Even after a year the ephemeris cannot differ very much from what is seen in the sky. For this reason Professor \textsc{Gauss} used set III of the elements above to calculate the ephemeris for the motion of Pallas in its next apparition, printed in the October issue of the \textit{Monatliche Correspondenz}.

\SourceLeaf{68}%
The second way in which Professor \textsc{Gauss} demonstrated the effect of the perturbations was to calculate purely elliptic elements that fit \textit{all} six oppositions \textit{as closely as possible} and nevertheless deviate considerably from the individual observed positions. We give this \textit{fourth} set of elements here as well:
\begin{center}\small
\newcommand{\GaussVIElementLeaders}{%
  .\hspace{.95em}.\hspace{.95em}.\hspace{.95em}.\hspace{.95em}.\hspace{.95em}%
  .\hspace{.95em}.\hspace{.95em}.\hspace{.95em}.\hspace{.95em}.\hspace{.95em}%
  .\hspace{.95em}.\hspace{.95em}.\hspace{.95em}.\hspace{.95em}.\hspace{.95em}.%
}
\begin{tabular}{cll}
$A$ & \GaussVIElementLeaders & \GaussVIAngle{221}{^0}{34}{'}{53}{''64}\\
$B$ & \GaussVIElementLeaders & \multicolumn{1}{r}{$770.5010$}\\
$C$ & \GaussVIElementLeaders & \GaussVIAngle{121}{}{8}{}{8}{.54}\\
$D$ & \GaussVIElementLeaders & \GaussVIAngle{172}{}{28}{}{12}{.43}\\
$E$ & \GaussVIElementLeaders & \GaussVIAngle{34}{}{37}{}{28}{.35}\\
$F$ & \GaussVIElementLeaders & $0.2447424$\\
$G$ & \GaussVIElementLeaders & $0.4422071$
\end{tabular}
\end{center}
\endgroup
The following table shows the errors of these elements:
\begin{center}\small
\setlength{\tabcolsep}{4pt}
\newcommand{\GaussVIError}[2]{\makebox[11pt][c]{$#1$}\makebox[39pt][r]{$#2$}}
\begin{tabular}{c|c|c}
Opposition & \multicolumn{2}{c}{Difference}\\
of & in heliocentric longitude & in geocentric latitude\\\hline
1803 & \GaussVIError{-}{111''00} & \GaussVIError{-}{8''31}\\
1804 & \GaussVIError{+}{59.18} & \GaussVIError{-}{36.67}\\
1805 & \GaussVIError{+}{19.92} & \GaussVIError{+}{0.07}\\
1807 & \GaussVIError{+}{85.77} & \GaussVIError{+}{25.01}\\
1808 & \GaussVIError{+}{135.88} & \GaussVIError{+}{28.72}\\
1809 & \GaussVIError{-}{216.54} & \GaussVIError{+}{83.01}
\end{tabular}\end{center}
The present paper is chiefly concerned with developing the methods by which these different sets of elements were found. In this respect it may be regarded as a kind of supplement to some sections of the \textit{Theoria motus corporum coelestium}. The problem of determining the orbit of a celestial body from four observed positions had already been discussed at length in that work, and two solutions were given, one for the case in which the orbit is still entirely unknown and one for the case in which the orbit, already known approximately, is to be improved. The latter could therefore also have been used to determine the three sets of elements given above. But in this particular case, where the four observations are oppositions, Professor \textsc{Gauss} found it more advantageous to use a quite different method, which the paper
\SourceLeaf{69}%
explains in detail and further illustrates by an example, the actual calculation of the third set of elements. Space here does not permit a closer description of this method.

The fourth set of elements was calculated according to the principles developed in Section 3 of Book 2 of the \textit{Theoria motus corporum coelestium}. The present paper also gives several additions in this respect that will, it is hoped, be welcome to astronomers. First comes a convenient way to calculate the differential changes in the heliocentric longitude and the geocentric latitude from the differential changes in the individual elements. Then comes a special procedure for determining the unknown quantities in accordance with the principle mentioned above. Let $w,w',w''$ and so on be the given linear functions of the unknown quantities $p,q,r$ and so on, and suppose that the aggregate $ww+w'w'+w''w''+$ and so on is to attain a minimum. It is then easy to obtain as many linear equations as there are unknowns, from which the unknowns must be determined by elimination. But when the number of unknowns is somewhat large, this elimination is extremely laborious because each equation contains all the unknowns. Professor \textsc{Gauss} has greatly shortened this work. Although he too reduces the solution to as many linear equations as there are unknowns, his equations are such that only the first contains all the unknowns. The second is free of $p$, the third of $p$ and $q$, the fourth of $p,q$ and $r$, and so on. Thus determining the unknowns in reverse order takes little effort. This method has the further advantage that one can state the least value of $ww+w'w'+w''w''+$ and so on in advance and use it as a check on the calculation by comparing it with the value computed afterwards, when the values found for the unknowns are substituted in $w,w',w''$, etc.
\par\bigskip\centerline{\rule{.24\textwidth}{.3pt}}
